% Generated draft: requires mathematical and source review.
\documentclass[11pt]{article}
\usepackage[margin=25mm]{geometry}
\usepackage{fontspec}
\setmainfont{Latin Modern Roman}
\usepackage{amsmath,amssymb,mathtools,unicode-math}
\setmathfont{Latin Modern Math}
\usepackage{enumitem,needspace,xurl,hyperref}
\setlength{\parindent}{0pt}
\setlength{\parskip}{5pt}
\setlength{\emergencystretch}{3em}
\newcommand{\N}{\mathbb{N}}
\newcommand{\Z}{\mathbb{Z}}
\newcommand{\Q}{\mathbb{Q}}
\newcommand{\R}{\mathbb{R}}
\newcommand{\C}{\mathbb{C}}
\newcommand{\F}{\mathbb{F}}
\newcommand{\A}{\mathbb{A}}
\newcommand{\PP}{\mathbb{P}}
\newcommand{\E}{\mathbb{E}}
\newcommand{\eps}{\varepsilon}
\newcommand{\dd}{\,\mathrm{d}}
\begin{document}
\section*{UnsolvedMath ID: 30001848}
\textbf{Source problem code:} catalogue entry
\subsection{Definitions and mathematical statement}
Fix an odd prime \(p\) and an integer \(g\geq 1\). The source does not explicitly define \(A_K\) or specify the admissible groups \(G\). The following formulation uses the conventional interpretation that \(A_K\) is the \(p\)-primary subgroup of the ideal class group of \(K\), and that \(G\) ranges over finite abelian \(p\)-groups of \(p\)-rank \(g\). These interpretations require source confirmation.

For an imaginary quadratic field \(K\), let \(\mathcal O_K\) be its ring of integers, let \(\operatorname{Cl}(\mathcal O_K)\) be its group of nonzero fractional ideals modulo principal fractional ideals, and set
\[
A_K=\{a\in\operatorname{Cl}(\mathcal O_K):p^n a=0\text{ for some integer }n\geq 0\}.
\]
For a finite abelian \(p\)-group \(H\), define
\[
r_p(H)=\dim_{\F_p}(H/pH).
\]
Interpret the \(p\)-rank of \(K\) as \(r_p(A_K)\). Let \(D_K<0\) denote the field discriminant. Interpret ordering by discriminant as ordering by \(|D_K|\), and count fields up to isomorphism. For real \(X>0\), define
\[
\mathcal K_{p,g}(X)=\{[K]:K\text{ is imaginary quadratic},\ |D_K|\leq X,\ r_p(A_K)=g\}.
\]
For a finite abelian \(p\)-group \(G\), let \(\operatorname{Aut}(G)\) be its group of automorphisms.

\textbf{Conjectural statement under these interpretations.} For every odd prime \(p\), every integer \(g\geq 1\), and every finite abelian \(p\)-group \(G\) with \(r_p(G)=g\), does the limiting frequency exist and satisfy
\[
\lim_{X\to\infty}
\frac{\#\{[K]\in\mathcal K_{p,g}(X):A_K\cong G\}}
{\#\mathcal K_{p,g}(X)}
=
\frac{p^{-g^2}}{|\operatorname{Aut}(G)|}
\prod_{k=1}^{g}\bigl(p^g-p^{g-k}\bigr)^2\,?
\]
The quotient is defined only when its denominator is nonzero. The source does not separately address the nonemptiness or infinitude of the conditioning family. The interpretation of discriminant ordering as ordering by absolute discriminant is also an explicit convention requiring confirmation.
\subsection{Short English statement}
Fix an odd prime \(p\) and a positive rank \(g\). Interpreting \(A_K\) as the \(p\)-primary ideal class group, the question asks whether imaginary quadratic fields of that rank, counted by absolute discriminant, have the displayed limiting distribution among finite abelian \(p\)-groups of rank \(g\). The group interpretation and counting conventions are explicit qualifications because the supplied statement does not fully specify them.
\subsection{Sources}
\begin{itemize}
\item Local registry: \texttt{data/unsolvedmath/problems.json}, permanent UnsolvedMath ID \(30001848\), title ``Fixed-Rank Cohen--Lenstra Distributions,'' source problem code \texttt{catalogue entry}, in the supplied Oberwolfach Reports Open Problems collection. No individual proposer is supplied; the proposed year is recorded as 2011.
\item Supplied report citation: \emph{Explicit Methods in Number Theory (2012)}, Oberwolfach Reports, pp. 1971--2044, DOI \texttt{10.4171/owr/2011/35}. Supplied literature-source label: ``original OWR report.'' \url{https://doi.org/10.4171/owr/2011/35}. The background describes Oberwolfach Reports as documenting mathematical workshops and their open-problem discussions. The citation's 2012 year and the record's 2011 proposed year are retained without reconciliation.
\item Supplied literature reference: ``Boston--Bush--Hajir p-class-tower heuristics,'' \url{https://doi.org/10.1007/s00208-016-1449-3}. The record does not provide a more detailed citation or identify a particular result supporting this formulation.
\item Supplied dated literature assessment, checked \texttt{2026-08-21}: the record labels the problem open and describes the full fixed-rank, arbitrary-odd-prime Cohen--Lenstra distribution for imaginary quadratic class groups as heuristic, with only partial cases proved. It lists the two references above as checked. This is supplied dated curation derived from the OWR record, not an independently verified literature assessment; its time-sensitive status and cited claims require independent verification.
\item The registry labels the statement \texttt{corrected\_verified} and reports that its mathematical content was checked against the official report, with notation, context, and question form normalized without changing the source-level claim. This is a supplied verification claim, not verification performed for this draft.
\end{itemize}
\end{document}
